\begin{table}[!htb]
\centering
\small
\setlength{\tabcolsep}{4pt}
\begin{tabular}{lcccccc}
\toprule
& \multicolumn{2}{c}{$g\to q\bar q$} & \multicolumn{2}{c}{$g\to gg$, $\nu=0$} & \multicolumn{2}{c}{$g\to gg$, $\nu=2$}\\
\cmidrule(lr){2-3}\cmidrule(lr){4-5}\cmidrule(lr){6-7}
$x$ & $C_F$ & $C_A$ & $C_A$ & $T_Fn_f$ & $C_A$ & $T_Fn_f$\\
\midrule
0.1 & $11.78598720$ & $-6.98106379$ & $5.75285579$ & $-4.56253144$ & $5.78524783$ & $-4.62731552$\\
0.3 & $3.56622290$ & $0.49139902$ & $5.07516361$ & $-4.70175009$ & $5.51057203$ & $-5.57256693$\\
0.5 & $2.13039556$ & $1.72312203$ & $4.89055459$ & $-4.74074074$ & $6.07573977$ & $-7.11111111$\\
\midrule
difference & $8.4\times10^{-8}$ & $4.1\times10^{-8}$ & $5.4\times10^{-8}$ & $3.5\times10^{-9}$ & $5.9\times10^{-8}$ & $1.8\times10^{-9}$\\
\bottomrule
\end{tabular}
\caption{The constant $r_0$ of the one-loop kernel from the regions (Sec.~\ref{app:regions}), coefficient by colour factor, for quark splittings ($gg\to ZZ$) and for the two components of gluon splittings ($gg\to HH$); by the symmetry $x\to1-x$, $x=0.7$ and $0.9$ give the same. The last row gives the largest difference between numerical region integrations and the analytic expressions of Eqs.~(\Lref{eq:r0g}) and~(\Lref{eq:r0q}) over the three values of $x$.}
\label{tab:r0}
\end{table}
